\section{Identidad estructural y magnitud de masa}
\label{vi:sec:introduccion}

Determinar la masa de una partícula exige precisar qué objeto persiste y
qué operación convierte su estructura en una magnitud. Este trabajo
construye ese enlace: de las extensiones compatibles y sus registros se
obtienen fibras de estados, y sobre ellas un operador positivo de masa.
Su especialización electrónica reúne dos lectores coincidentes en una
misma evaluación escalar. La tesis es que la identidad, la ecuación de
masa y el valor evaluado pertenecen a una cadena matemática común; la
comparación física examina después la realización de esa cadena en una
coordenatización de medida. Así, los argumentos y coeficientes de la ecuación
conservan una procedencia estructural anterior al contraste.

La dinámica de trayectoria incorpora la historia en esa magnitud: determina
el soporte celular exacto de la trayectoria central, reconstruye dos memorias
que una estadística de conteo no distingue, diagonaliza el lector de retorno
mediante caracteres finitos y caracteriza la conservación de masa por el
conmutador entre el operador de masa y un transporte TPK. Estas cuatro
operaciones actúan sobre el mismo estado y no reemplazan el atlas, la ecuación
electrónica ni sus definiciones y pruebas previas.

Se introduce un núcleo formal matemático denominado Holografía Modular
Triádica, en adelante HMT, compuesto por tres objetos. \emph{All Possible
Paths}, en adelante APP, es la estructura aritmética de caminos sobre
\((\mathbb Z/9\mathbb Z)^2\). La adyacencia viene dada por el grafo de
Cayley aditivo con generadores cardinales \(\pm(1,0)\) y \(\pm(0,1)\);
la realización celular periódica de ese soporte es toroidal. Sus
operaciones de acumulación aditiva y composición
multiplicativa se distinguen de las evaluaciones que determinan sus
coordenadas. TRIT es la unidad de firma orientada \(\{-1,0,+1\}\):
determina el régimen local y la orientación de la operación. El
\emph{Topological Phase Kernel}, en adelante TPK, constituye la dinámica
del núcleo; compone selección, transporte, actualización de memoria y
retorno sobre el estado enriquecido. La Mecánica Dimensional desarrolla
las realizaciones de esa estructura en líneas de magnitud y espacios de
estados físicos. La exposición construye estos objetos desde sus reglas,
antes de utilizar las coordenadas internas obtenidas en los operadores de masa.

El desarrollo obtiene resultados definidos en cada uno de esos niveles. La
composición de \(108\) actualizaciones del TPK produce un estado terminal
que conserva el libro incidencial, sus dos canales transversales y la firma
dodecafásica; la transformada orbital reconstruye de manera reversible el registro
\(K\). Sobre esta genealogía se construyen el atlas de
\(81/56/13/324\), el operador positivo de masa de cada fibra, la excepción
escalar del electrón central y las reglas de composición anteriores a la
evaluación. El contraste posterior mantiene separados valores centrales,
intervalos, límites, masas de polo, masas corrientes y anchuras. La aportación
consiste, por tanto, en una cadena efectiva que enlaza el objeto persistente,
su memoria de ruta, el operador que actúa sobre su fibra y el tipo preciso de
la magnitud evaluada.

En ese mismo orden causal, el censo completo
\(104\,976\to468\to243\subset729\) precede a toda selección regional. La
bijección \(\beta:\mathbb F_3^6\to E_9^3\) identifica el ambiente de
\(729\) posiciones, mientras la diferencia de umbral produce la corona
\(271=243+27+1\) y la completación exacta
\(1000=729+243+27+1\). La medida uniforme de la completación es compatible
con las proyecciones dimensionales. Estas operaciones construyen el ambiente
y su borde; la selección posterior conserva su dominio propio.

\subsection{La estructura que precede a la cifra}

La masa no es la única información de una partícula. También intervienen
el soporte de persistencia, sus orientaciones, la conjugación, las
representaciones y las reglas de composición. En el planteamiento HMT,
estos datos no se agregan al final a una cifra aislada: pertenecen a la
construcción sobre la que actúa el funcional de masa. Dos historias pueden
compartir una lectura escalar y conservar información diferente en sus
hojas o en su memoria. Por ello, la igualdad de valores debe distinguirse
de la igualdad de estructuras.

La distinción es particularmente clara en la composición. Una tabla puede
sumar coordenadas ya evaluadas; una dinámica debe resolver primero si dos
historias pueden componer su fase, su borde, sus orientaciones y sus
predecesores. Los registros se componen en ese nivel. La firma se extrae
después, y el carácter evalúa esa firma. La diferencia entre esas
operaciones aporta el lugar matemático para estudiar estados compuestos,
enlaces y resonancias.

La pregunta por qué vale lo que vale una magnitud queda así vinculada a
otra más precisa: qué elementos de la construcción determinan su valor y
qué información continúa disponible después de evaluarlo. El artículo
desarrolla esa pregunta desde las reglas internas del corpus. La
comparación física posterior examina sus realizaciones sin intervenir en
la selección de los coeficientes que se someten a contraste.

\subsection{Jerarquía de la exposición}

El núcleo formal se reproduce de manera común con los artículos
precedentes de la serie. La repetición cumple una función de autonomía:
cada texto debe poder leerse sin disponer de los demás. Las ampliaciones
específicas del presente trabajo aparecen después del núcleo y se
identifican por los objetos que construyen.

Las coordenadas generadas, la constante de estructura fina y la constante
reducida de Planck intervienen cuando las ecuaciones másicas las
requieren. El par angular precede a los canales y a sus momentos; los
registros preceden a la firma; las fibras preceden a los operadores que
actúan en ellas. El electrón central se conserva como construcción
estructural propia, aun cuando su magnitud se exprese en una escala
cronológica compartida.

Dentro de esta jerarquía, selección, transporte y actualización actúan
sobre un mismo estado enriquecido. La composición de \(108\) pasos
produce el estado terminal; su proyección al libro dodecafásico, la
evaluación de los canales \(90/120\), la extracción firmada
\(U_{12}^{\rm sgn}\), la transformada orbital y la reconstrucción
integral establecen la cadena
\[
 X_{\rm seed}\longmapsto X_{\rm term}\longmapsto
 C_{\rm term}\longmapsto U_{\rm term}\longmapsto K.
\]
El registro, sus \(12\) componentes, sus cocientes y su memoria proceden
de esa evolución. Las coordenatizaciones de canales, firma y período conservan en un
mismo dominio la hoja, la orientación, el acarreo y la frontera requeridos
por la evaluación de masa. La composición~\eqref{vi:dep:cadena-terminal-completa}
reúne el productor y su reconstrucción, con las pruebas de actualización e
inversión contiguas; la identidad~\eqref{vi:dep:K-reversible} precisa la
transformación reversible y su subred integral.

La definición de partícula prepara la clasificación. A continuación se
desarrollan la conjugación y la holonomía, el atlas, los observables
internos y la evaluación de masa. La distribución por colecciones utiliza
esa arquitectura, de manera que cada sector se presenta con su soporte,
su operación y su lectura. Kepler--Sommerfeld se sitúa en el desarrollo
orbital de la acción, donde pueden especificarse los mapas y las
condiciones que justifican la relación; no se introduce como una
semejanza histórica sin contenido operativo.

\subsection{La comparación como resultado de una cadena declarada}

El contraste requiere más que dos columnas de números. Debe saberse qué
se ha evaluado, en qué unidad y bajo qué definición del observable. Una
masa de polo, una masa corriente a una escala declarada y un límite
experimental no son intercambiables. Una anchura expresa información
dinámica distinta de la posición de una masa. El volumen del inventario
es útil cuando conserva estas distinciones, porque permite examinar
colecciones completas y no sólo ejemplos favorables.

Por esa razón, la presentación reúne cada salida con su procedencia y
con las condiciones de su realización. Una identidad interna, la
evaluación de una firma especificada y la identificación de esa firma
con una especie son afirmaciones relacionadas que poseen demostraciones
propias. Exponerlas consecutivamente permite mostrar qué determina el
modelo en cada nivel y cuál es el contenido concreto del contraste.

La unidad argumental del artículo es, en definitiva, la continuidad entre
persistencia, clasificación, operación y medida. La masa es una salida
de esa continuidad; la estructura que la produce permanece disponible
para comprender sus relaciones con las demás lecturas de la partícula.

% BEGIN SINTESIS_ADITIVA_20260922
\par
La sección~\ref{delta22:vi:composicion} incorpora la longitud estructural
antes del reloj: la incidencia marcada produce \(N=5760\),
\(r_N=5759/23040\) y
\(L_\alpha=\alpha_{\rm HMT}^{16}r_NL_\star\)
en~\eqref{delta22:vi:Lalpha}; sólo después se fijan \(t_0\),
\(R_{108}=L_\alpha\), \(Q_L\), \(m_{\rm clk}\) y
\(b_e(Q_L)\). La ley de cotransformación
\eqref{delta22:vi:be-cotransformacion} muestra por qué el factor
electrónico no puede congelarse al variar la sección longitudinal o la
acción. Las relaciones operatorias \eqref{delta22:vi:R01} y
\eqref{delta22:vi:R02} eliminan la longitud de acción y determinan el
coeficiente gravitatorio único~\eqref{delta22:vi:G-unica}; la fórmula
circular~\eqref{delta22:vi:G-circular} se obtiene después. La masa, su
carácter de ruta y sus tablas permanecen causalmente anteriores a esa
identificación radial.
% END SINTESIS_ADITIVA_20260922
